\subsection{\problemblock{*Inertial} Data Block}\label{sec:block-inertial}

TAOS maintains an inertial-platform coordinate system that is fixed in inertial
space. The Euler angles \statevar{yawi}, \statevar{pitchi}, and
\statevar{rolli} are referenced to this system, and position, velocity, and
acceleration can be output in it as described in \cref{sec:inertial-platform}.

The platform is initially aligned with the local geodetic horizon coordinate system
and located at the trajectory initial position. Its $x$ axis points north, its $y$ axis
points east, and its $z$ axis points down toward the earth. This alignment and
position can be changed at the beginning of any segment with an
\problemblock{*inertial} block.

\begin{table}[htbp]
  \centering
  \caption{Inertial-platform alignment coordinate systems.}
  \label{tab:inertial-platform-alignments}
  \begin{tabularx}{0.82\linewidth}{@{}>{\ttfamily}lX@{}}
    \toprule
    \normalfont Name & Description \\
    \midrule
    body       & Body-axis system \\
    ecfc       & Earth-centered, earth-fixed system; optional east and down vectors \\
    geocentric & Local geocentric horizon; optional longitude/latitude/radius position \\
    geodetic   & Local geodetic horizon; optional longitude/latitude/altitude position \\
    velocity   & $x$ parallel to velocity, with $z$ coplanar with the geodetic $z$ axis \\
    wind       & Wind-axis system \\
    \bottomrule
  \end{tabularx}
\end{table}

The default alignment is geodetic. All coordinate systems in
\cref{tab:inertial-platform-alignments} are defined in \cref{sec:coordinate-systems}.

\taossubhead{Geocentric and Geodetic Alignment}

The data block

\begin{lstlisting}[style=taosmanual]
*inertial platform alignment geocentric
\end{lstlisting}

aligns the inertial system with the current local geocentric horizon and locates its
origin at the vehicle center of mass. The words \texttt{platform} and
\texttt{alignment} are optional; the shorter form \texttt{*inertial geocentric}
is equivalent.

The platform may instead be placed at a specified longitude, latitude, and altitude,
and at a specified time:

\begin{lstlisting}[style=taosmanual]
*inertial platform lat=28.8 long=-81.5 alt=550 time=100
\end{lstlisting}

This aligns with the geodetic horizon at $28.8^\circ$~N, $81.5^\circ$~W, an
altitude of 550~ft, and a mission time of 100~s. Alignment time is required to account
for earth rotation. For a geodetic alignment, \statevar{long}, \statevar{lat}, and
\statevar{alt} specify the location; altitude is optional and defaults to zero, or sea
level. For geocentric alignment, \statevar{long}, \statevar{lat}, and
\statevar{rcm} are used.

\taossubhead{Body, Wind, and Velocity Alignment}

The inertial platform can be aligned with the body axes with

\begin{lstlisting}[style=taosmanual]
*inertial platform body
\end{lstlisting}

Body-attitude information is not available at the beginning of the first trajectory
segment because the guidance scheme has not yet been executed. The platform should
therefore not be aligned with the body or wind systems on the first segment. An
initial dummy segment with a final condition of \texttt{tseg>0} can first calculate
the attitude; at the start of the second segment the platform can then be aligned with
the body, wind, or velocity system.

For body, wind, or velocity alignment, the platform is located at the vehicle center of
mass. An alignment time may still be supplied. This aligns the platform as though
the current position and velocity existed at the specified mission time.

\taossubhead{ECFC Alignment}

For earth-centered, earth-fixed alignment, the platform origin is the vehicle center of
mass and its axes are parallel to the ECFC unit vectors. The orientation may be
changed by entering east and down vectors. Only two unit vectors are needed; TAOS
computes the remaining vector to form a right-handed system. Components are given
in ECFC coordinates with \statevar{eastx}, \statevar{easty},
\statevar{eastz}, \statevar{downx}, \statevar{downy}, and
\statevar{downz}. If only one vector is supplied, the default is used for the other.
For example,

\begin{lstlisting}[style=taosmanual]
*inertial ecfc eastx=0.7071 easty=0.7071 eastz=0.0
               downx=0.0   downy=0.0    downz=-1
\end{lstlisting}

twists the east and north inertial vectors by approximately $45^\circ$.
